% One-time imported draft. Editable native source; no Markdown regeneration.
\section{Step 8 Separate vertical magnitude from
direction}\label{007-bresenham-line-s08-step-8-separate-vertical-magnitude-from-direction}

Begin with direction +1. If the signed y difference is negative, negate
it and choose direction \ensuremath{-}1. \texttt{bgez} has an explicit \texttt{nop}
delay slot. The result is a nonnegative vertical magnitude plus a
separate signed step.

\subsection{Write or
inspect}\label{007-bresenham-line-s08-write-or-inspect}

Append the absolute-value and direction block.

\MARSCodeLabel{MIPS assembly}
\begin{minted}{mars-idea}
    li $s6, 1
    bgez $s5, absolute_ready
    nop
    subu $s5, $zero, $s5
    li $s6, -1
absolute_ready:
\end{minted}

\subsection{Expected
observation}\label{007-bresenham-line-s08-expected-observation}

\texttt{\$s5} is the vertical magnitude D; \texttt{\$s6} is +1 or \ensuremath{-}1.
For the default they are three and one.

\subsection{Check your
understanding}\label{007-bresenham-line-s08-check-your-understanding}

Why retain a separate y direction?

\begin{itemize}
\tightlist
\item
  \textbf{A} The magnitude must keep a negative sign.
\item
  \textbf{B} The distance magnitude alone cannot say whether y increases
  or decreases.
\item
  \textbf{C} All negative slopes are out of bounds.
\end{itemize}
